\documentclass[UTF8,11pt]{ctexart}
\usepackage[a4paper,margin=24mm]{geometry}
\usepackage{amsmath,amssymb,mathtools}
\usepackage{xcolor}
\usepackage{hyperref}
\usepackage{enumitem}
\hypersetup{colorlinks=true,linkcolor=blue!55!black,urlcolor=blue!55!black}
\setlist{nosep}

\newcommand{\Bool}{\ensuremath{\mathbb{B}}}
\newcommand{\Clean}{\ensuremath{\mathsf{Clean}}}
\newcommand{\InputChanged}{\ensuremath{\mathsf{InputChanged}}}
\newcommand{\NodeChanged}{\ensuremath{\mathsf{NodeChanged}}}
\newcommand{\compat}{\ensuremath{\operatorname{compatible}}}
\newcommand{\mat}{\ensuremath{\operatorname{M}}}

\title{CSS 增量匹配器：四种变体、Hoare Logic 与 Kani 验证}
\author{css-bitvector-compiler}
\date{\today}

\begin{document}
\maketitle

\begin{abstract}
本文给出 bit、tri、recursive-tri 与 quad 四种增量 CSS selector
匹配器的统一语义。四种算法共享同一 selector 状态机、DOM arena 和 mutation
模型，只在缓存抽象与失效处理上不同。本文使用 Hoare triple 描述关键优化的
正确性条件，并说明 Kani model checking harness 所证明的局部转换定律。
\end{abstract}

\section{共享状态机}

令编译后的状态集合为 $Q$。每个 DOM 节点 $n$ 接收 parent 与 previous-sibling
两个输入通道，并产生物化输出：
\[
P_n,S_n,O_n\in\Bool^{|Q|},\qquad \Bool=\{0,1\}.
\]

例如 selector \texttt{.a .b > .c} 产生状态 $q_0,q_1,q_2$：
\[
\begin{aligned}
q_0(n)&=p_{\mathtt{.a}}(n),\\
q_1(n)&=P_n[q_0]\land p_{\mathtt{.b}}(n),\\
q_2(n)&=P_n[q_1]\land p_{\mathtt{.c}}(n),
\qquad q_2\in Q_{\mathrm{accept}}.
\end{aligned}
\]
descendant combinator 对中间状态增加传播：
\[
O_n[q_i]=q_i(n)\lor P_n[q_i].
\]
child 不传播；adjacent sibling 从 $S_n$ 读取 predecessor state。

节点 dirty 状态构成全序：
\[
\Clean<\InputChanged<\NodeChanged.
\]
\NodeChanged 表示本地事实变化，不能被输入缓存跳过；\InputChanged 表示仅
$P_n$ 或 $S_n$ 变化。

\section{Bit：物化输出缓存}

Bit 缓存 $O_n$，对 \NodeChanged 和 \InputChanged 均完整执行节点程序 $F_n$：
\[
\mathsf{step}_{\mathrm{bit}}(n)=
\begin{cases}
F_n(P_n,S_n),&d_n\in\{\NodeChanged,\InputChanged\},\\
O_n,&d_n=\Clean.
\end{cases}
\]
只有当新旧 $O_n$ 不同时，children 的 parent input 与下一个 sibling 的
sibling input 才会被标记为失效。

\section{Tri：输入依赖抽象}

Tri 为每个输入 bit 保存 requirement：
\[
\mathcal R=\{\bot,0,1\},
\]
其中 $\bot$ 表示本次执行未读取。具体化关系为：
\[
\gamma(\bot)=\{0,1\},\qquad
\gamma(0)=\{0\},\qquad
\gamma(1)=\{1\}.
\]
因此：
\[
\compat(R,I')\iff\forall i,\ I'[i]\in\gamma(R[i]).
\]

Tri skip 的 Hoare triple 是：
\[
\left\{
\begin{array}{l}
O_n=F_n(I)\\
\land\ R\text{ 记录了执行 }F_n(I)\text{ 时所有实际读取}\\
\land\ \compat(R,I')
\end{array}
\right\}
\quad\mathsf{skip}\quad
\left\{F_n(I')=O_n\right\}.
\]

该规则要求严格保持短路求值。例如
$p_{\mathtt{.b}}(n)\land P_n[q]$ 在左侧为 false 时不能读取 $P_n[q]$，
否则会把无关输入错误记录为依赖并损失优化。

\section{Recursive-tri：调度前检查}

Tri 在访问节点后决定 skip；recursive-tri 在父输出变化时先检查 child：
\[
P_c\ne P'_c\Longrightarrow
\begin{cases}
\text{不调度 }c,&\compat(R_c,P'_c),\\
d_c\gets\InputChanged,&\text{otherwise}.
\end{cases}
\]

其 Hoare triple 为：
\[
\begin{aligned}
&\left\{d_c=\Clean\land\compat(R_c,P'_c)\right\}\\
&\qquad\mathsf{do\_not\_schedule}(c)\\
&\left\{O'_c=O_c\land
\neg\operatorname{invalidate}(\operatorname{desc}(c))\right\}.
\end{aligned}
\]

当前实现采用逐层 admission check。child 不被调度意味着其输出不变，所以失效
不会继续进入 descendants。显式 folded subtree requirement 是可选的未来优化，
不是当前正确性的前提。

\section{Quad：输出函数抽象}

Quad 的抽象域为：
\[
\mathcal Q=\{\mathbf0,\mathbf1\}
\cup\{\pi^P_i\mid 0\le i<|Q|\}
\cup\{\pi^S_i\mid 0\le i<|Q|\}.
\]
物化函数为：
\[
\begin{aligned}
\mat(\mathbf0,P,S)&=0,&\mat(\mathbf1,P,S)&=1,\\
\mat(\pi^P_i,P,S)&=P[i],&\mat(\pi^S_i,P,S)&=S[i].
\end{aligned}
\]

计算分两阶段：先求 local/predecessor transition $r$，再组合 descendant
传播 $r\lor\pi^P_i$。由于 $\mathcal Q$ 不能表示任意 OR，按当前输入专门化：
\[
\operatorname{specializeOr}(r,c,I)=
\begin{cases}
(r,r=1),&\mat(r,I)=1,\\
(c,r=0),&\mat(r,I)=0.
\end{cases}
\]

令结果和 decision 为 $(q,d)$，则 OR refinement 的 Hoare triple 是：
\[
\left\{(q,d)=\operatorname{specializeOr}(r,c,I)
\land\operatorname{holds}(d,I')\right\}
\quad x\gets\mat(q,I')\quad
\left\{x=\mat(r,I')\lor\mat(c,I')\right\}.
\]

accept state 必须具体化。若
$(b,d)=\operatorname{specializeConcrete}(q,I)$，则：
\[
\left\{\operatorname{holds}(d,I')\right\}
\quad x\gets\mat(b,I')\quad
\left\{x=\mat(q,I')\right\}.
\]

decision 失效时执行完整重算；decision 仍成立时只重新 materialize abstract
output。这样纯传播节点无需重新执行 selector predicates。

\section{Mutation 与结构不变量}

\begin{itemize}
  \item class、id、attribute 和派生 pseudo 变化属于 \NodeChanged。
  \item \texttt{:hover} 向下派生，\texttt{:focus-within} 向上派生。
  \item sibling 插入/删除使后续节点的 nth position 失效。
  \item accept bit 只表示当前节点匹配，不能从祖先直接传播。
  \item attribute shortcut 仅在局部输出与匹配集合均不变时停止传播。
\end{itemize}

\section{Kani proof harness}

proof harness 直接调用生产代码中的纯函数。Kani 固定 parent/sibling vector
长度为 2，但 index、所有输入 bit 和 Quad tag 都是符号值。

\begin{enumerate}
  \item \nolinkurl{tri_requirement_reuse_is_sound}：compatible 的 Zero/One requirement 保证被读取 bit 不变。
  \item \nolinkurl{quad_projection_materializes_current_input}：parent/sibling projection 精确读取对应输入。
  \item \nolinkurl{specialized_or_refines_boolean_or}：decision 成立时专门化结果等价于原布尔 OR。
  \item \nolinkurl{concrete_accept_reuse_is_sound}：decision 成立时 concrete accept 可安全复用。
\end{enumerate}

运行：
\begin{verbatim}
cargo xtask verify
cargo xtask check
cargo xtask corpus
\end{verbatim}

形式化证明覆盖逐 bit 的局部转换定律。完整 parser、无界 DOM 和递归树结构尚未
被归纳证明；端到端组合由 13 个真实站点与独立 naive oracle 做 differential
validation。

\section{当前限制}

当前不支持 \texttt{:has}、general sibling、\texttt{:last-*}、
\texttt{:nth-last-*} 和 selector-list pseudo。无界 DOM 的归纳证明适合后续用
Verus、Lean 或 Coq 建模。

\end{document}
